% !TeX program = pdflatex
% The sign pattern of Sun's trigonometric permanents and a counterexample
% to his Conjecture 4.6.  Build: pdflatex note.tex (run twice). See BUILD.md.
\documentclass[10pt]{article}

\usepackage[T1]{fontenc}
\usepackage{lmodern}
\usepackage[margin=1in]{geometry}
\usepackage{amsmath,amssymb,amsthm}
\usepackage{booktabs}
\usepackage{microtype}
\usepackage[hidelinks]{hyperref}
\usepackage{xurl}

\newtheorem{theorem}{Theorem}[section]
\newtheorem{lemma}[theorem]{Lemma}
\theoremstyle{remark}
\newtheorem{remark}[theorem]{Remark}

\DeclareMathOperator{\per}{per}
\newcommand{\Z}{\mathbb{Z}}
\newcommand{\Q}{\mathbb{Q}}
\newcommand{\F}{\mathbb{F}}

\title{The sign pattern of Sun's trigonometric permanents\\
and a counterexample to his Conjecture 4.6}
\author{John Erlbacher\thanks{Independent Researcher. Email: \texttt{erlbacher.research@gmail.com}. ORCID: 0009-0003-6851-4139.}}
\date{June 2026}

\begin{document}
\maketitle

\begin{abstract}
In \emph{Arithmetic properties of some permanents} (arXiv:2108.07723),
Zhi-Wei Sun proved that
$s_n := (2^{(n-1)/2}/\sqrt n)\per[\sin 2\pi jk/n]_{1\le j,k\le(n-1)/2}$ and
$s'_p := (\sqrt p/2^{(p-1)/2})\per[\csc 2\pi jk/p]_{1\le j,k\le(p-1)/2}$ are
integers, and conjectured (Conjecture~4.6) that (i) $n \mid s_n$ for odd
composite $n$, and (ii) for odd primes, $s_p<0 \iff p\equiv 5 \pmod{12}$ and
$s'_p<0 \iff p\equiv 7\pmod 8$. We refute part~(ii) --- hence the conjecture
as stated --- at $p=29$, the first prime beyond Sun's published table:
$s_{29}=1\,053\,859>0$ although $29\equiv 5\pmod{12}$, and
$s'_{29}=-4\,806\,838\,304<0$ although $29\equiv 5\pmod 8$. The mod-$12$
clause fails again at $p=41$ and $53$, the mod-$8$ clause at $p=61$. Every
value is certified exactly, by at least two of four independently written
exact programs, via a cyclotomic identity evaluated in two ways:
big-integer arithmetic in $\Z[x]/(x^{4n}-1)$ (the counterexample pair,
among others), and finite-field specialization with CRT uniqueness
proofs. The pipeline reproduces all nineteen of Sun's
published values, and every new prime value satisfies his proven
congruences.
Part~(i) survives all tests (the fifteen odd composites $n\le 65$) and
remains open.
\end{abstract}

\section{Introduction}\label{sec:intro}

In \cite{Sun2021}, Sun studies permanents of trigonometric matrices and
proves, among many results, the following integrality theorem
(Theorem~1.6 of \cite{Sun2021}). For odd $n>1$ write $m=(n-1)/2$. Then
\begin{equation}\label{eq:defs}
s_n := \frac{2^{m}}{\sqrt n}\,
       \per\Bigl[\sin 2\pi\frac{jk}{n}\Bigr]_{1\le j,k\le m} \in \Z
\qquad\text{and}\qquad
s'_p := \frac{\sqrt p}{2^{m}}\,
       \per\Bigl[\csc 2\pi\frac{jk}{p}\Bigr]_{1\le j,k\le m} \in \Z,
\end{equation}
the first for all odd $n>1$, the second for all odd primes
$p$.\footnote{In the arXiv \TeX{} source the second definition is typeset
``$s_p'=:$'', an evident typo for ``$:=$''; the intended reading is
unambiguous from the surrounding theorem.} Here $\per A = \sum_{\tau\in
S_m}\prod_{j=1}^{m} a_{j\tau(j)}$ is the permanent. The same theorem proves
the congruences
\begin{equation}\label{eq:cong}
s_p \equiv (-1)^{(p+1)/2} \pmod p
\qquad\text{and}\qquad
s'_p \equiv 1 \pmod p
\end{equation}
for every odd prime $p$. In Remark~1.6, Sun reports the values
$s_3,\dots,s_{23}$ and $s'_3,\dots,s'_{23}$ (nineteen integers, reproduced
in Table~\ref{tab:calibration} below), and in Section~4 he poses, motivated
by those data:\footnote{Conjectures in \cite{Sun2021} are numbered by
section; the statement quoted here is the sixth conjecture of Section~4,
i.e.\ Conjecture~4.6 of the compiled paper (verified against the compiled
PDF of arXiv:2108.07723v7). It should not be confused with the identically
numbered conjectures of other papers of Sun.}

\begin{quote}
\textbf{Conjecture 4.6} (Sun \cite{Sun2021}). \emph{(i) If $n>1$ is odd and
composite, then $s_n\equiv 0\pmod n$. (ii) Let $p$ be an odd prime. Then
\[
s_p<0 \iff p\equiv 5 \pmod{12},
\qquad\text{and}\qquad
s'_p<0 \iff p\equiv 7 \pmod 8.
\]}
\end{quote}

\noindent
We disprove the conjecture by computing $s_p$ and $s'_p$ exactly for the
nine primes $29\le p\le 61$ and $s_n$ for all fifteen odd composite
$n\le 65$. Both clauses of part~(ii) already fail at $p=29$, the first
prime beyond the published table (Theorem~\ref{thm:main}). Part~(i), by
contrast, is \emph{not} refuted: $n\mid s_n$ holds in every tested case
(Remark~\ref{rem:parti}), and that clause remains an open and, on this
evidence, plausible conjecture.

The matrices are dense, with $m$ up to $32$, and the values grow to $31$
digits, so the computation's trustworthiness deserves comment. Every value
reported here was computed by
\emph{exact} methods only: an identity in the cyclotomic ring
$\Z[\zeta_{4n}]$ (Lemma~\ref{lem:cyclo}) reduces each $s_n$, $s'_p$ to an
exactly verifiable integer equation, which was evaluated in two ways ---
exact big-integer arithmetic in $\Z[x]/(x^{4n}-1)$ (for the
counterexample pair and the small cases), and finite-field
specialization at several primes $q\equiv 1\pmod{4n}$ with
Chinese-remainder uniqueness proofs (Section~\ref{sec:method}). Five
independently written programs were used --- four exact, plus a
high-precision check by Glynn's formula; every value reported here was
certified by at least two of the exact programs, with no disagreement
observed anywhere, and every new prime value satisfies the proven
congruences~\eqref{eq:cong} (Section~\ref{sec:verification}).

Related recent work concerns the \emph{determinant} analogues: Gao and Guo
\cite{GaoGuo2025} resolve a conjecture of Sun on a trigonometric
$\sec$-determinant (the $\csc$ analogue having been settled in earlier
work of the same authors, cited there). We are not aware of any prior work on the permanent
conjecture treated here; as of June~2026 the arXiv record of \cite{Sun2021}
is unchanged since 2022 (v7), and neither of the sequences $(s_n)$,
$(s'_p)$ appears in the OEIS, so Tables~\ref{tab:primes}
and~\ref{tab:composites} also extend the published data --- thirty new
exact values --- beyond Sun's table.

\paragraph{Methods and AI-use disclosure.}
The counterexamples were found and verified by Demonstrandum, a
verification-first multi-agent AI search-and-reasoning pipeline built on
Claude language models (Anthropic):
agents froze the conjecture statement from the arXiv \TeX{} source of
\cite{Sun2021}, wrote the attack code, and then wrote the verification
layers of Section~\ref{sec:verification} independently of one another and
of the attack code. An adversarial audit was performed by a model from a
different family (Codex, GPT-5.5, OpenAI), which probed the statement
reading and the verification mathematics and independently recomputed the
kill pair $(s_{29},s'_{29})$ and $s_{41}$ at 100--120-digit precision. All
accept paths run on exact integer and finite-field arithmetic; floating
point appears only in advisory cross-checks and in the generously padded
magnitude bounds of one CRT program (Section~\ref{sec:method}), and every
value is also certified by a layer whose bounds are exact-rational or
that needs no magnitude bound at all.
The text of this note was likewise drafted with the Claude-based pipeline;
the author directed the work and takes responsibility for its content.
Every number appearing in this note was recomputed or re-verified from the
certificate data during its preparation (June 2026); see
Section~\ref{sec:verification}.

\section{The counterexamples}\label{sec:results}

\begin{theorem}\label{thm:main}
Conjecture~4.6 of \cite{Sun2021} is false. Both clauses of part~(ii) fail
at $p=29$, where $29\equiv 5\pmod{12}$ and $29\equiv 5\pmod 8$:
\[
s_{29} \;=\; 1\,053\,859 \;>\; 0,
\qquad\qquad
s'_{29} \;=\; -4\,806\,838\,304 \;<\; 0,
\]
so the implication $p\equiv 5\pmod{12}\Rightarrow s_p<0$ fails, and so does
the implication $s'_p<0\Rightarrow p\equiv 7\pmod 8$. Moreover the
first clause fails at every prime $p\equiv 5\pmod{12}$ in the computed
range ($p=29,41,53$, with $s_{41}=5\,574\,476\,521$ and
$s_{53}=1\,540\,679\,755\,916\,971$, both positive), and the second clause
fails again at $p=61\equiv 5\pmod 8$, where
\[
s'_{61}\;=\;-2\,904\,784\,276\,786\,469\,053\,142\,518\,062\,479\;<\;0.
\]
\end{theorem}

\begin{proof}
A finite exact computation. The values of $s_p$ and $s'_p$ for the nine
primes $29\le p\le 61$ are listed in Table~\ref{tab:primes}; each is
certified, via the cyclotomic method of Section~\ref{sec:method}, by at
least two independently written exact programs with disjoint prime sets
(the kill pair additionally by the exact big-integer implementation and a
third CRT program), validated as described in
Section~\ref{sec:verification}. The residues are elementary. Since both
violations are sign statements about nonzero integers, no boundary or
strictness issue arises; note that $s'_p=0$ is in any case excluded by the
proven congruence~\eqref{eq:cong}.
\end{proof}

\begin{table}[ht]
\centering\footnotesize
\setlength{\tabcolsep}{3.5pt}
\begin{tabular}{@{}rrrrrl@{}}
\toprule
$p$ & $p\bmod{12}$ & $p\bmod 8$ & $s_p$ & $s'_p$ & violates \\
\midrule
29 & 5 & 5 & $1\,053\,859$ & $-4\,806\,838\,304$ & both clauses \\
31 & 7 & 7 & $10\,542\,977$ & $-1\,518\,806\,869\,720$ & --- \\
37 & 1 & 5 & $1\,519\,663\,259$ & $43\,041\,655\,439\,377$ & --- \\
41 & 5 & 1 & $5\,574\,476\,521$ & $88\,188\,655\,594\,502\,880$ & mod-12 clause \\
43 & 7 & 3 & $453\,435\,884\,081$ & $7\,817\,331\,967\,711\,147\,274$ & --- \\
47 & 11 & 7 & $4\,570\,760\,060\,257$ & $-208\,210\,243\,110\,377\,949\,730$ & --- \\
53 & 5 & 5 & $1\,540\,679\,755\,916\,971$ & $1\,364\,915\,376\,262\,405\,923\,148\,317$ & mod-12 clause \\
59 & 11 & 3 & $493\,044\,351\,638\,203\,633$ & $7\,313\,054\,784\,852\,201\,235\,037\,895\,089\,487$ & --- \\
61 & 1 & 5 & $3\,864\,901\,746\,617\,921\,299$ & $-2\,904\,784\,276\,786\,469\,053\,142\,518\,062\,479$ & mod-8 clause \\
\bottomrule
\end{tabular}
\caption{Exact values of $s_p$ and $s'_p$ for the nine primes beyond Sun's
published table. The mod-12 clause of Conjecture~4.6(ii) is violated iff
($p\equiv 5\pmod{12}$ and $s_p>0$) or ($p\not\equiv 5$ and $s_p<0$); the
mod-8 clause iff ($p\equiv 7\pmod 8$ and $s'_p>0$) or ($p\not\equiv 7$ and
$s'_p<0$). Every value satisfies the proven
congruences~\eqref{eq:cong}.}
\label{tab:primes}
\end{table}

\begin{remark}[Part (i) is not refuted]\label{rem:parti}
The divisibility clause was tested exactly for \emph{all} fifteen odd
composite $n\le 65$, including the prime squares $9$, $25$, $49$ and the
values $25,35,49,55,65$ coprime to~$3$; Table~\ref{tab:composites} lists $s_n$ and
the integer quotient $s_n/n$. The divisibility $n\mid s_n$ holds in every
case. Part~(i) of Conjecture~4.6 therefore survives this attack and remains
open; our refutation of the conjecture (a conjunction) rests entirely on
part~(ii).
\end{remark}

\begin{table}[ht]
\centering\footnotesize
\setlength{\tabcolsep}{3.5pt}
\begin{tabular}{@{}rrr@{\hspace{2em}}rrr@{}}
\toprule
$n$ & $s_n$ & $s_n/n$ & $n$ & $s_n$ & $s_n/n$ \\
\midrule
 9 & $9$ & $1$                              & 45 & $696\,741\,294\,375$ & $15\,483\,139\,875$ \\
15 & $45$ & $3$                             & 49 & $3\,078\,195\,713\,761$ & $62\,820\,320\,689$ \\
21 & $2\,835$ & $135$                       & 51 & $113\,109\,498\,706\,113$ & $2\,217\,833\,307\,963$ \\
25 & $63\,125$ & $2\,525$                   & 55 & $9\,553\,828\,212\,328\,125$ & $173\,705\,967\,496\,875$ \\
27 & $59\,049$ & $2\,187$                   & 57 & $89\,237\,136\,634\,668\,369$ & $1\,565\,563\,800\,608\,217$ \\
33 & $-1\,643\,895$ & $-49\,815$            & 63 & $14\,357\,464\,732\,984\,700\,313$ & $227\,896\,265\,602\,931\,751$ \\
35 & $334\,186\,125$ & $9\,548\,175$        & 65 & $23\,595\,726\,326\,056\,828\,125$ & $363\,011\,174\,247\,028\,125$ \\
39 & $9\,154\,298\,673$ & $234\,725\,607$   &    &  & \\
\bottomrule
\end{tabular}
\caption{The part-(i) test set: exact $s_n$ for all odd composite
$n\le 65$. The quotient $s_n/n$ is an integer in every case, so part~(i) of
Conjecture~4.6 is \emph{not} refuted. (The values for $n=9,15,21$ agree
with Sun's table.)}
\label{tab:composites}
\end{table}

\begin{remark}[Why the pattern broke]\label{rem:pattern}
Each sign law of part~(ii) was fitted to eight primes containing only two
negatives each ($s_p<0$ at $p=5,17$; $s'_p<0$ at $p=7,23$). In the extended
range the first law collapses completely: $s_p>0$ for \emph{every} prime
$29\le p\le 61$, so the only known negative values of $s_p$ remain $p=5$
and $17$, and the law fails at every tested $p\equiv 5\pmod{12}$ beyond the
published table. Negative values of $s'_p$ do persist --- at $p=31$ and
$47$, both $\equiv 7\pmod 8$, the second law happens to conform --- but
they also appear at $p=29$ and $61$, both $\equiv 5\pmod 8$. Galois
equivariance does not separate the permanent from its determinant
cousins here: the automorphism $\sigma_a\colon\zeta_p\mapsto\zeta_p^a$ of
$\Q(\zeta_p)$ multiplies $\per M$ (the matrix of Lemma~\ref{lem:cyclo}
with $n=p$; its entries $\zeta_p^{jk}-\zeta_p^{-jk}$ lie in
$\Z[\zeta_p]$) by the Legendre symbol $(\tfrac ap)$
--- the computation at the heart of Sun's integrality proof --- so the
permanent, too, is a Galois eigenvector, and that symmetry constrains
the sign of neither quantity. What the determinant analogues studied by Sun and in
\cite{GaoGuo2025} additionally possess is a character factorization:
their determinants factor over Dirichlet characters into character sums
--- in \cite{GaoGuo2025}, into special values of Dirichlet $L$-functions
--- and structure of that kind can force residue-class behaviour. No
analogous factorization is known for the permanent, and we know of no
mechanism that could enforce a residue-class sign rule for $s_p$ or
$s'_p$. We do not propose a replacement law: the observed signs describe
the computed range $p\le 61$ only.
\end{remark}

\section{The cyclotomic verification method}\label{sec:method}

The permanents in~\eqref{eq:defs} are sums of $m!$ products of irrational
reals, the results are integers as large as $10^{30}$, and a sign is
precisely what is at stake; the following identity removes all analysis
from the problem.

For odd $n>1$ let $\zeta=e^{2\pi i/(4n)}$, a primitive $4n$-th root of
unity, so that $i=\zeta^{n}$ and $e^{2\pi i\,jk/n}=\zeta^{4jk}$, and let
\[
G \;=\; \sum_{a=0}^{n-1} \zeta^{4a^2} \;\in\; \Z[\zeta]
\]
be the quadratic Gauss sum. By Gauss's theorem (see, e.g.,
\cite[Thm.~1.5.2]{BEW1998}), $G=\sqrt n$ if $n\equiv 1\pmod 4$ and
$G=i\sqrt n$ if $n\equiv 3\pmod 4$. Set $\widetilde G=G$ in the first case
and $\widetilde G=\zeta^{-n}G=-iG$ in the second, so that
$\widetilde G=\sqrt n$ for every odd $n$; in particular
$\sqrt n\in\Z[\zeta]$.

\begin{lemma}\label{lem:cyclo}
Let $n>1$ be odd, $m=(n-1)/2$, and $M_{jk}=\zeta^{4jk}-\zeta^{-4jk}$ for
$1\le j,k\le m$. Then
\begin{equation}\label{eq:sin-id}
\per M \;=\; s_n\, i^{m}\sqrt n \qquad\text{in } \Z[\zeta].
\end{equation}
If moreover $n=p$ is prime, the matrix $U$ with entries
$U_{jk}=p\,(\zeta^{4jk}-\zeta^{-4jk})^{-1}$ has entries in $\Z[\zeta]$, and
\begin{equation}\label{eq:csc-id}
\per U \;=\; s'_p\, p^{\,m-1}\sqrt p\; i^{-m} \qquad\text{in } \Z[\zeta].
\end{equation}
\end{lemma}

\begin{proof}
From $\sin\theta=(e^{i\theta}-e^{-i\theta})/(2i)$ we get $\sin(2\pi
jk/n)=M_{jk}/(2i)$, and multilinearity of the permanent in the rows gives
$\per[\sin 2\pi jk/n]=(2i)^{-m}\per M$; substituting into~\eqref{eq:defs}
yields~\eqref{eq:sin-id}. For~\eqref{eq:csc-id}, $\csc(2\pi
jk/p)=2i\,M_{jk}^{-1}$, so $\per[\csc 2\pi jk/p]=(2i)^{m}p^{-m}\per U$, and
\eqref{eq:defs} gives $\per U=s'_p\,p^{m}i^{-m}/\sqrt p
=s'_p\,p^{m-1}\sqrt p\,i^{-m}$. The entries of $U$ are algebraic integers:
$\zeta^{4jk}-\zeta^{-4jk}=\zeta^{-4jk}(\omega^{jk}-1)$ with
$\omega=\zeta^{8}$ a primitive $p$-th root of unity and $p\nmid jk$, and
$\omega^{jk}-1$ divides $p=\prod_{t=1}^{p-1}(1-\omega^{t})$ in
$\Z[\omega]$.
\end{proof}

Both sides of \eqref{eq:sin-id} and \eqref{eq:csc-id} lie in $\Z[\zeta]
\cong \Z[x]/(\Phi_{4n}(x))$, where $\Phi_{4n}$ is the $4n$-th cyclotomic
polynomial, and the right-hand sides are nonzero integer multiples of a
fixed ring element; hence each identity determines the integer $s_n$
(resp.\ $s'_p$) uniquely and reduces its computation to finite integer
arithmetic. We used two independent implementations.

\paragraph{Implementation A: exact arithmetic in $\Z[x]/(x^{4n}-1)$.}
Represent elements of $\Z[\zeta]$ by integer polynomials modulo
$x^{4n}-1$, with $x$ standing for $\zeta$. The permanent of $M$ (or $U$) is
computed by Ryser's inclusion--exclusion formula \cite{Ryser1963} over this
ring, a Gray-code walk over the $2^m$ column subsets; for speed, ring
elements are packed into pairs of nonnegative big integers by Kronecker
substitution $x\mapsto 2^{K}$, with $K$ chosen from a rigorous a-priori
$L^1$ coefficient bound so that no base-$2^{K}$ digit can overflow, and
cyclic reduction modulo $x^{4n}-1$ becomes arithmetic modulo $2^{4nK}-1$
(every decode is re-encoded and asserted equal as an integrity check). The
program then \emph{solves} for the unique integer $t$ with $P\equiv
t\,B\pmod{\Phi_{4n}(x)}$ over $\Z$, where $P$ is the computed permanent and
$B=x^{nm}\,\widetilde G(x)$ (resp.\ the right side of~\eqref{eq:csc-id});
$\Phi_{4n}$
is computed by exact division of $x^{4n}-1$ by lower cyclotomic
polynomials. For the $\csc$ case the entries $U_{jk}$ are obtained by the
extended Euclidean algorithm over $\Q[x]$ modulo $\Phi_{4p}$ and verified
by exact back-multiplication $U_{jk}\cdot(x^{4jk}-x^{-4jk})\equiv
p\pmod{\Phi_{4p}}$ over~$\Z$. No floating point and no finite fields are
used anywhere. On a desktop machine this implementation delivers $s_{29}$
in about four and a half minutes and $s'_{29}$ in about seven.

\paragraph{Implementation B: finite-field specialization and CRT
uniqueness.}
For a prime $q\equiv 1\pmod{4n}$, the field $\F_q$ contains elements $w$ of
exact multiplicative order $4n$; any such $w$ is a root of $\Phi_{4n}$
modulo $q$, so reduction modulo $q$ and evaluation $x\mapsto w$ is a ring
homomorphism $\Z[x]/(\Phi_{4n}(x))\to\F_q$. The images of
\eqref{eq:sin-id}--\eqref{eq:csc-id} are then identities in $\F_q$ in which
the image $g$ of $\widetilde G$ is computed from the defining sum $\sum_a
w^{4a^2}$, multiplied by $w^{-n}$ when $n\equiv 3\pmod 4$ (so no
square-root sign ambiguity can arise; the implementations assert
$g^2\equiv n \pmod q$), the permanent of the specialized matrix is computed
modulo $q$, and $s_n \bmod q$ is read off. Repeating this for several
random primes $q$ and combining by the Chinese remainder theorem with
$\prod q > 4B$, where $B$ is a bound on the magnitude of the
target integer ($|s_n|\le 2^m m!$ suffices for the $\sin$ case; for the
$\csc$ case an explicit bound follows from the chord inequality $|\sin \pi
x|\ge 2\,\mathrm{dist}(x,\Z)$), determines the integer \emph{uniquely}
--- a proof, not a probabilistic check, provided $B$ itself is rigorous.
Three independently written programs of this type were run, with disjoint
random prime sets and different permanent kernels: a Gray-code Ryser in C
with Montgomery arithmetic (59-bit primes; generously padded
floating-point bounds, with every value additionally confirmed at a
held-out prime not used in the reconstruction), a plain 128-bit-integer
Ryser in C (fresh primes from $[2^{60},2^{61})$; bounds evaluated in
exact rational arithmetic), and a pure-Python subset-dynamic-programming
permanent (57-bit primes; exact rational bounds and a held-out prime per
value), the last using the recursion
$f(S)=\sum_{j\in S}A_{|S|,j}\,f(S\setminus\{j\})$
rather than Ryser's formula. The $[2^{60},2^{61})$ run certified 32 values
--- the kill pair, every $s_n$ for odd $31\le n\le 65$, every $s'_p$ for
$31\le p\le 61$, and four values from Sun's table --- in about 22 minutes
of wall-clock time.

Implementations A and B share no code; A was written by an agent with no
access to the discovery code. In addition to the exact layers, two
floating-point recomputations were performed as advisory cross-checks:
Glynn's formula \cite{Glynn2010} --- a third permanent algorithm --- in
60-digit arithmetic, and, during the adversarial audit, independent
100--120-digit recomputations of $s_{29}$, $s'_{29}$ and $s_{41}$ by both
Glynn's and Ryser's formulas directly from the real entries (residuals
below $10^{-86}$ in all cases).

\section{Verification and calibration}\label{sec:verification}

\paragraph{Calibration on Sun's table.}
Before computing anything new, the pipeline was calibrated against the
nineteen values published in Remark~1.6 of \cite{Sun2021}
(Table~\ref{tab:calibration}): three of the five programs ---
implementation A, the first CRT program, and the subset-DP program ---
reproduced all nineteen exactly, and the remaining two were first
validated on subsets containing negative entries (the clean-room CRT
program on $s_{17}$, $s_{23}$, $s'_{17}$, $s'_{23}$; the Glynn check on
its anchors). The four negative entries ($s_5$, $s_{17}$, $s'_7$,
$s'_{23}$) pin the sign conventions end-to-end: a global sign error in any
layer --- in the Gauss sum, in the power of $i$, or in the permanent
kernel --- would flip at least one of them.

\begin{table}[ht]
\centering\footnotesize
\setlength{\tabcolsep}{5.5pt}
\begin{tabular}{@{}lrrrrrrrrrrr@{}}
\toprule
$n$    & 3 & 5 & 7 & 9 & 11 & 13 & 15 & 17 & 19 & 21 & 23 \\
$s_n$  & $1$ & $-1$ & $1$ & $9$ & $1$ & $51$ & $45$ & $-239$ & $913$ & $2\,835$ & $12\,145$ \\
\midrule
$p$    & 3 & 5 & 7 & 11 & 13 & 17 & 19 & 23 & & & \\
$s'_p$ & $1$ & $1$ & $-6$ & $111$ & $261$ & $6\,784$ & $245\,101$ & $-7\,094\,142$ & & & \\
\bottomrule
\end{tabular}
\caption{Calibration: the nineteen values published by Sun
\cite[Remark~1.6]{Sun2021}, all reproduced exactly --- by three of the
five programs in full, by the other two on validation subsets --- before
any new value was computed.}
\label{tab:calibration}
\end{table}

\paragraph{Cross-checks on the new values.}
The exact programs of Section~\ref{sec:method} agree wherever their
coverage overlaps, and every entry of Tables~\ref{tab:primes}
and~\ref{tab:composites} is certified by at least two of
them.\footnote{Coverage detail: implementation A computed the kill pair
and the small composites exactly; the first CRT program computed every
value; the clean-room CRT program certified the 32 values listed above;
the subset-DP program covers the kill pair and $s_{41}$. Thus the headline
values are certified by three or four exact programs plus the
floating-point layers. The supplementary values $s_{53}$ and $s'_{61}$
rest on the two independent CRT programs alone and were not part of the
100-digit audit recomputation.} Beyond
mutual agreement: every $s_p$, $s'_p$ for $29\le p\le 61$ satisfies the
congruences~\eqref{eq:cong}, which are \emph{proven} in \cite{Sun2021} ---
the only theorem available to test the new values against (for the kill
pair: $s_{29}\equiv 28\equiv(-1)^{15}$ and $s'_{29}\equiv 1\pmod{29}$); the
first and third CRT programs confirmed each of their values --- between
them, every value in this note --- at a held-out prime; and the adversarial
audit (three rounds, Codex/GPT-5.5) probed the statement provenance
(permanent vs.\ determinant, the half-range index set, the conjecture
numbering) and recomputed the headline values independently at high
precision.

\paragraph{Reproducibility.}
All computations were re-run while preparing this note (June 2026). A
reader can reproduce the kill in under two minutes: a one-page
floating-point script (Ryser's formula in IEEE doubles) reproduces
$s_{17}$, $s'_7$, $s'_{23}$ and the kill pair to the nearest integer in
about a second, and the pure-Python
exact program (implementation B, subset-DP kernel) re-derives all nineteen
calibration values, the kill pair, and $s_{41}$, with full CRT uniqueness
proofs, in under a minute. While drafting this note we additionally
re-verified the congruences~\eqref{eq:cong} and the residues and quotients
of Tables~\ref{tab:primes}--\ref{tab:composites}, and recomputed, by a
freshly written floating-point Gray-code Ryser, every table entry with
$m\le 18$ (all nineteen calibration values; $s_p$, $s'_p$ for
$p=29,31,37$; $s_n$ for $n=25,27,33,35$).

\paragraph{Data availability.}
The certificate bundle (all five programs, their logs, machine-readable
certificates with the 33 exact integers of
Tables~\ref{tab:primes}--\ref{tab:composites} and their provenance, and
the frozen arXiv source of \cite{Sun2021}) will be made available at
\url{https://github.com/demonstrandum-research/[REPO]}.

\begin{thebibliography}{9}\small\setlength{\itemsep}{1pt}

\bibitem{BEW1998}
B.~C. Berndt, R.~J. Evans, and K.~S. Williams,
\emph{Gauss and Jacobi Sums},
Canadian Mathematical Society Series of Monographs and Advanced Texts
\textbf{21}, Wiley, New York, 1998.

\bibitem{GaoGuo2025}
L.~Gao and X.~Guo,
\emph{Integrality of a trigonometric determinant arising from a conjecture
of Sun},
preprint, arXiv:2512.24012 (2025).

\bibitem{Glynn2010}
D.~G. Glynn,
\emph{The permanent of a square matrix},
European J. Combin.\ \textbf{31} (2010), 1887--1891.

\bibitem{Ryser1963}
H.~J. Ryser,
\emph{Combinatorial Mathematics},
Carus Mathematical Monographs \textbf{14},
Mathematical Association of America, 1963.

\bibitem{Sun2021}
Z.-W. Sun,
\emph{Arithmetic properties of some permanents},
preprint, arXiv:2108.07723v7 (2021--2022).

\end{thebibliography}

\end{document}
